\chapter*{Abstract}
\label{chap:abstract}
The Method of Loci is a memorisation technique in which information is associated with locations along a familiar route. Virtual Reality (VR) can make this process visible and interactive, allowing users to place mnemonic cues in a virtual environment. However, a predefined environment may have little personal meaning to the person using it. A virtual version of a familiar home could provide a more recognisable structure, while also allowing the user to adjust the space to suit the task.

In this thesis, we present a VR memory palace system that supports the creation and furnishing of personalised environments and the placement of visual mnemonic cues. The current study compares a generic palace with a personalised palace based on the participant's home. A within-subjects design is used to examine recall, cognitive load, enjoyment, perceived competence, effort, and presence. The protocol developed through pilot sessions, with changes to word-list difficulty, list length, and the presence questionnaire.

The analysis snapshot available on 22 September 2026 contains 18 complete immediate-recall pairs, 16 pairs at the nominal 24-hour follow-up, and 14 pairs at the nominal one-week follow-up. Immediate mean recall accuracy is 72.39\% in the generic condition and 78.67\% in the personalised condition. The estimated paired difference is 6.28 percentage points, with a 95\% confidence interval from $-2.51$ to 15.06 percentage points. These values are provisional: the analysis workbook flags six blocking data-quality errors and 130 unreviewed, uncounted responses, and the dataset includes multiple protocol versions. Consequently, this draft does not establish a recall benefit or equivalence between conditions.

The present contribution is the design of the personalised palace workflow and a documented evaluation framework. Content personalisation and auditory mnemonic generation remain part of the broader proposal, but are outside the comparison reported here. The thesis identifies the data reconciliation and experimental controls needed before drawing final conclusions about spatial personalisation.

